\section{Conclusion}

We presented a controlled empirical study of how data augmentation affects the double descent phenomenon in neural networks. Our main finding is unambiguous: \textbf{ResNets without augmentation show a clear double descent on CIFAR-10}, with test error rising 16\% between the 697K-parameter trough and the 2.78M-parameter peak, while \textbf{all four augmentation strategies completely eliminate this peak}, yielding monotonically decreasing error curves with 65--69\% lower error at large model sizes. Three mechanisms appear to jointly drive this effect: effective sample size increase, loss landscape smoothing, and inductive bias regularization.

A key negative finding is that \textbf{MixUp is insensitive to its interpolation coefficient $\alpha$}: values from 0.1 to 1.0 produce nearly identical results, suggesting the existence---rather than the intensity---of label mixing is what matters. We also find that \textbf{MLPs show no double descent} under any condition, implicating convolutional or residual architecture components as necessary conditions for the phenomenon in our setting.

Practically, our results support a surprisingly simple recommendation: \emph{use any standard data augmentation}. Flip+Crop alone is sufficient to eliminate double descent on CIFAR-10---you do not need MixUp or Cutout for this purpose, though they do give better final accuracy. If we had to pick one takeaway for practitioners, it would be that the choice of augmentation strategy matters much less for avoiding double descent than for final accuracy: even the cheapest geometric augmentation appears to be enough.

The most important open question raised by this work is whether the same holds at scale. Our experiments top out at 2.78M parameters on 32$\times$32 images; it is entirely possible that at the parameter counts used in modern large models, the interpolation threshold reappears even under augmentation. We would not assume our findings transfer without verification.
